
\documentclass{article}
\usepackage{colortbl}
\usepackage{makecell}
\usepackage{multirow}
\usepackage{supertabular}

\begin{document}

\newcounter{utterance}

\centering \large Interaction Transcript for game `ifeval', experiment `ifeval', episode 240 with Qwen3.5{-}4B.
\vspace{24pt}

{ \footnotesize  \setcounter{utterance}{1}
\setlength{\tabcolsep}{0pt}
\begin{supertabular}{c@{$\;$}|p{.15\linewidth}@{}p{.15\linewidth}p{.15\linewidth}p{.15\linewidth}p{.15\linewidth}p{.15\linewidth}}
    \# & \multicolumn{2}{c}{Player} && \multicolumn{2}{c}{Game Master} \\
    \hline

    \theutterance \stepcounter{utterance}  
    & & & \multicolumn{4}{p{0.6\linewidth}}{
        \cellcolor[rgb]{0.9,0.9,0.9}{
            \makecell[{{p{\linewidth}}}]{
                \texttt{\tiny{[P1$\langle$GM]}}
                \texttt{Write a casual blog post about how the outer solar system is different from the inner solar system, and what that means for the possibility of life. Wrap your entire response with double quotation marks. Your response should contain 17 or more sentences.} \\
            }
        }
    }
     \\ \\

    \theutterance \stepcounter{utterance}  
    & \multicolumn{4}{p{0.6\linewidth}}{
        \cellcolor[rgb]{0.561,0.737,0.561}{
            \makecell[{{p{\linewidth}}}]{
                \texttt{\tiny{[P1$\rangle$GM]}}
                \texttt{"Have you ever looked up at the night sky and wondered why the planets seem so different from one another? The inner solar system is a hot, rocky playground dominated by Mercury, Venus, Earth, and Mars. These worlds are close to the sun, which means they are scorching and dry with little to no atmosphere to keep them cool. In contrast, the outer solar system is a frozen, icy frontier filled with gas giants and distant dwarf planets. Jupiter, Saturn, Uranus, and Neptune are massive balls of hydrogen and helium that spin incredibly fast. The temperatures down there are so low that water exists as ice rather than liquid or steam. This huge difference in temperature creates a completely different environment for any potential life to thrive. Life as we know it needs liquid water, which is surprisingly rare in the inner system but common in the outer system. The icy moons of Jupiter and Saturn, like Europa and Enceladus, actually have subsurface oceans hidden beneath thick layers of ice. These hidden oceans might be warmer due to tidal heating caused by the massive gravity of their parent planets. While the inner planets are great for studying geology and meteorology, the outer planets offer a better chance for microbial life. The pressure and radiation in the outer system are extreme, but the deep oceans could be shielded from these harsh conditions. Scientists are currently sending rovers to Mars to look for ancient signs of life in the inner system. At the same time, they are planning missions to drill into the ice of Europa to sample that hidden water. It is fascinating to think that the very conditions that make Earth habitable might be found far away in the cold dark. The outer solar system reminds us that life could be hiding in places we never expected to look. Exploring these distant worlds is one of the most exciting frontiers in modern astronomy today. We might just find our next best friend or even a new home among the stars soon."} \\
            }
        }
    }
    & & \\ \\

    \theutterance \stepcounter{utterance}  
    & & & \multicolumn{2}{p{0.3\linewidth}}{
        \cellcolor[rgb]{0.9,0.9,0.9}{
            \makecell[{{p{\linewidth}}}]{
                \texttt{\tiny{[GM$|$GM]}}
                {'startend:quotation': {}, 'length_constraints:number_sentences': {'relation': 'at least', 'num_sentences': 17}}
            }
        }
    }
    & & \\ \\

    \theutterance \stepcounter{utterance}  
    & & & \multicolumn{2}{p{0.3\linewidth}}{
        \cellcolor[rgb]{0.9,0.9,0.9}{
            \makecell[{{p{\linewidth}}}]{
                \texttt{\tiny{[GM$|$GM]}}
                \texttt{game\_result = WIN} \\
            }
        }
    }
    & & \\ \\

\end{supertabular}
}

\end{document}
